\emph{સાબિતી-પ્રણાલીઓ} વાક્યો વચ્ચેના તાર્કિક પરિણામ અને સુસંગતતાને
લક્ષણબદ્ધ કરવાની કેવળ વાક્યરચનાત્મક પદ્ધતિઓ આપે છે. \emph{સ્વાભાવિક નિગમન}
એવી એક પ્રણાલી છે. તેમાં \emph{નિષ્પત્તિ} સૂત્રોના વૃક્ષરૂપે હોય છે.
વૃક્ષની ટોચે આવેલાં સૂત્રો \emph{ધારણાઓ} છે. નિષ્પત્તિ સાચી હોય તે માટે
બાકીનાં બધાં સૂત્રોને કોઈ ને કોઈ \emph{અનુમાનનિયમ}થી યોગ્ય ઠેરવવાં પડે.
દરેક સંયોજક અને પરિમાણક માટે એક પ્રવેશ અને એક લોપનિયમ છે. દાખલા તરીકે,
$!A$ને $\Elim{\lif}$ વડે યોગ્ય ઠેરવીએ તો તેનાં પૂર્વગામી સૂત્રો
(\emph{આધારવાક્યો}) કોઈ $!B$ માટે $!B \lif !A$ અને $!B$ હોવાં જોઈએ.
કેટલાક નિયમો ધારણાઓને \emph{નિવૃત્ત} કરવાની પણ છૂટ આપે છે. દાખલા તરીકે,
$\Intro{\lif}$ વડે $!B$ પરથી $!A \lif !B$ મેળવીએ તો $!B$ સુધીની
નિષ્પત્તિમાં ધારણા તરીકે આવેલા $!A$ના દાખલાઓને નિવૃત્ત કરી શકાય છે. એ
દાખલાઓને અપાયેલી સંજ્ઞા અનુમાન પાસે પણ નોંધાય છે.

છેલ્લું સૂત્ર $!A$ હોય અને બધી ધારણાઓ નિવૃત્ત થઈ ગઈ હોય એવી નિષ્પત્તિ
હોય તો $!A$ \emph{પ્રમેય} છે. આપણે $\Proves !A$ લખીએ છીએ. બધી અનિવૃત્ત
ધારણાઓ કોઈ ગણ $\Gamma$માં હોય તો $!A$ \emph{$\Gamma$માંથી નિષ્પન્ન કરી
શકાય છે} અને $\Gamma \Proves !A$ લખીએ છીએ. જો $\Gamma \Proves \lfalse$
તો $\Gamma$ \emph{અસંગત} છે, અન્યથા \emph{સુસંગત}. આ ખ્યાલો પરસ્પર
સંબંધિત છે. દાખલા તરીકે, $\Gamma \Proves !A$ ત્યારે અને માત્ર ત્યારે જ
જ્યારે $\Gamma \cup \{\lnot !A\}$ અસંગત હોય. અનુરૂપ અર્થવિષયક ખ્યાલો સાથે
પણ સંબંધ છે. દાખલા તરીકે, $\Gamma \Proves !A$ હોય તો $\Gamma \Entails !A$.
સાબિતી-પ્રણાલીનો આ ગુણધર્મ કે $\Gamma$માંથી નિષ્પન્ન કરી શકાય તે બધું
$\Gamma$નું તાર્કિક પરિણામ હોય છે, તેને \emph{યથાર્થતા} કહે છે. નિષ્પત્તિની
લંબાઈ પર આગમનથી \emph{યથાર્થતાનું પ્રમેય} સાબિત થાય છે. તેમાં બતાવીએ છીએ
કે દરેક અનુમાનના આધારવાક્યો પોતપોતાની અનિવૃત્ત ધારણાઓનાં તાર્કિક પરિણામ હોય
તો તેનું તારણ પણ પોતાની અનિવૃત્ત ધારણાઓનું તાર્કિક પરિણામ રહે છે.
